\question[2]
Decodiere die Bitfolge mit dem angegebenen Huffman-Baum. Gib für jedes Zeichen die Codierung an.

\texttt{101010001100111100110111011110001}

Hinweis zur Darstellung des Huffman-Baums: Der Baum ergibt sich durch
Drehen der abgebildeten Zeilen um 90 Grad nach rechts. Die Punkte
geben die Ebenen des Baumes an. Knoten, die nur Zahlen aber keine
Zeichen enthalten, sind mit einem \# markiert.
\begin{lstlisting}
...2 r
..4 #
...2 u
.7 #
...2 g
..3 #
...1 k
12 #
..3 e
.5 #
...1 z
..2 #
...1 w
\end{lstlisting}
\begin{solutionbox}{2cm}
\begin{lstlisting}
g->101, e->01, w->000, u->110, r->111, z->001, k->100
Kodiertes Wort: gewuerzgurke
\end{lstlisting}
\end{solutionbox}
